% For hyperlinked PDF, suitable for viewing on a computer, use this:
\documentclass[letterpaper,12pt,titlepage,oneside,final]{book}

% For PDF suitable for double-sided printing, use this:
%\documentclass[letterpaper,12pt,titlepage,openright,twoside,final]{book}

\newcommand{\package}[1]{\textbf{#1}} % package names in bold text
\newcommand{\cmmd}[1]{\textbackslash\texttt{#1}} % command name in tt font
\newcommand{\href}[1]{#1}

% This package allows if-then-else control structures.
\usepackage{xifthen}
\newboolean{PrintVersion}
\setboolean{PrintVersion}{false}

\usepackage{amsmath,amssymb,amstext}
\usepackage[pdftex]{graphicx}
\usepackage[dvipsnames,hyperref]{xcolor}
\usepackage{microtype}
\usepackage[round]{natbib}
\usepackage{rotating}
\usepackage{multicol}
\usepackage{booktabs}
\usepackage{tipa}
\let\ipa\textipa

% Set figure directories
\graphicspath{{../figures/}}

% Custom hyphenation
\hyphenation{re-arranged}

% Includes a figure. Usage:
%
% \fig{filename}{width}{caption}{shortcaption}
%
% - Filename will be used as label
% - Width is in proportion of column width
% - Shortcaption is optional
\newcommand{\fig}[4]{
  \begin{figure}[ht!]
    \centering
    \includegraphics[width=#2\columnwidth]{#1}
    \ifthenelse{\isempty{#4}}{\caption{#3}}{\caption[#3]{#4}}
    \label{fig:#1}
  \end{figure}}

% Font stuff
\usepackage{inconsolata}
\usepackage{fourier}
\usepackage[T1]{fontenc}

% argmin and argmax aren't normal?
\DeclareMathOperator*{\argmin}{arg\,min}
\DeclareMathOperator*{\argmax}{arg\,max}
\newcommand*{\bigoh}{O}
% typing mathbf over and over again is garbage
\newcommand*{\V}[1]{\mathbf{#1}}%
% binding
\newcommand*{\bind}{\circledast}
% We got room; big summations please!
\everymath{\displaystyle}

% For fancy boxes
\usepackage{tikz}
\usetikzlibrary{shapes,decorations}
\tikzstyle{roundbox}=[draw=black!80, very thick,
    rectangle, rounded corners, inner sep=10pt, inner ysep=20pt]
\tikzstyle{fancytitle}=[fill=black!80, text=white]
\newcommand{\probbox}[1]{%
  \begin{center}
    \begin{tikzpicture}
      \node [roundbox] (box){%
        \begin{minipage}{0.8\textwidth}
          #1
        \end{minipage}
      };
      \node[fancytitle, right=10pt, rounded corners] at (box.north west) {%
        \textbf{Problem statement}
      };
    \end{tikzpicture}
\end{center}}
\newcommand{\inoutbox}[2]{%
  \begin{center}
    \begin{tikzpicture}
      \node [roundbox] (box){%
        \begin{minipage}{0.9\textwidth}
          \begin{multicols}{2}
            \begin{center}
              #1
              \vfill
              \columnbreak

              #2
            \end{center}
          \end{multicols}
        \end{minipage}
      };
      \node[fancytitle, right=8.4em, rounded corners] at (box.north west) {%
        \textbf{Inputs}
      };
      \node[fancytitle, left=8.2em, rounded corners] at (box.north east) {%
        \textbf{Outputs}
      };
      \draw[very thick] (box.north) -- (box.south);
    \end{tikzpicture}
\end{center}}

% Fix error using \copyright with T1 fontenc
\renewcommand*\copyright{{\usefont{OT1}{lmr}{m}{n}\textcopyright}}

% Do this last
\usepackage[pdftex,pagebackref=false]{hyperref}
\hypersetup{
    plainpages=false,
    unicode=false,
    pdftoolbar=true,
    pdfmenubar=true,
    pdffitwindow=false,
    pdfstartview={FitH},
    pdftitle={Dynamical systems in spiking neuromorphic hardware},
    pdfauthor={Aaron R. Voelker},
    pdfsubject={Dynamical Neuromorphic Systems},
    pdfkeywords={nef} {nengo} {spiking neural networks} {dynamical systems}
                {computational neuroscience} {neuromorphics},
    pdfnewwindow=true,
    colorlinks=true,
    linkcolor=BrickRed,
    citecolor=PineGreen,
    urlcolor=cyan
}
\ifthenelse{\boolean{PrintVersion}}{
  \hypersetup{
    citecolor=black,
    filecolor=black,
    linkcolor=black,
    urlcolor=black}
}{}

\setlength{\marginparwidth}{0pt} % width of margin notes
% N.B. If margin notes are used, you must adjust \textwidth, \marginparwidth
% and \marginparsep so that the space left between the margin notes and page
% edge is less than 15 mm (0.6 in.)
\setlength{\marginparsep}{0pt} % width of space between body text and margin notes
\setlength{\evensidemargin}{0.125in} % Adds 1/8 in. to binding side of all
% even-numbered pages when the "twoside" printing option is selected
\setlength{\oddsidemargin}{0.125in} % Adds 1/8 in. to the left of all pages
% when "oneside" printing is selected, and to the left of all odd-numbered
% pages when "twoside" printing is selected
\setlength{\textwidth}{6.375in} % assuming US letter paper (8.5 in. x 11 in.) and
% side margins as above
\raggedbottom

% The following statement specifies the amount of space between
% paragraphs. Other reasonable specifications are \bigskipamount and \smallskipamount.
\setlength{\parskip}{\medskipamount}

% The following statement controls the line spacing.  The default
% spacing corresponds to good typographic conventions and only slight
% changes (e.g., perhaps "1.2"), if any, should be made.
\renewcommand{\baselinestretch}{1}

% Force each section of the front pages to start on a recto page.
% Also ensure a page number is not printed on an otherwise blank verso page.
\let\origdoublepage\cleardoublepage
\newcommand{\clearemptydoublepage}{%
  \clearpage{\pagestyle{empty}\origdoublepage}}
\let\cleardoublepage\clearemptydoublepage

\begin{document}

% suppress the page number and headers/footers.
\pagestyle{empty}
\pagenumbering{roman}

\begin{titlepage}
  \begin{center}
    \vspace*{1.0cm}

    \Huge
    {\bf Dynamical Systems in Spiking Neuromorphic Hardware}

    \vspace*{1.0cm}

    \normalsize
    by \\

    \vspace*{1.0cm}

    \Large
    Aaron Russell Voelker \\

    \vspace*{3.0cm}

    \normalsize
    A thesis \\
    presented to the University of Waterloo \\
    in fulfillment of the \\
    thesis requirement for the degree of \\
    Doctor of Philosophy \\
    in \\
    Computer Science \\

    \vspace*{2.0cm}

    Waterloo, Ontario, Canada, 2019 \\

    \vspace*{1.0cm}

    \copyright\ Aaron R. Voelker 2019 \\
  \end{center}
\end{titlepage}

% no headers, but yes page numbers (starting from ii.)
\pagestyle{plain}
\setcounter{page}{2}

\cleardoublepage

\phantomsection
\addcontentsline{toc}{chapter}{Examining Committee Membership}
\begin{center}\textbf{Examining Committee Membership}\end{center}

\noindent
The following served on the Examining Committee for this thesis.
The decision of the Examining Committee is by majority vote.

\noindent\begin{tabular}{@{}ll}\\
    External Examiner & Adrienne Fairhall \\ & Professor \\ \\
    Supervisor & Chris Eliasmith \\ & Professor \\ \\
    Internal Member & Jeff Orchard \\ & Associate Professor \\ \\
    Internal Member & Ali Ghodsi \\ & Professor \\ \\
    Internal-external Member & Sue Ann Campbell \\ & Professor
\end{tabular}

\cleardoublepage
\phantomsection
\addcontentsline{toc}{chapter}{Author's Declaration}

\noindent
This thesis consists of material all of which I authored or co-authored: see Statement of Contributions included in the thesis. This is a true copy of the thesis, including any required final revisions, as accepted by my examiners.

\bigskip

\noindent
I understand that my thesis may be made electronically available to
the public.

\cleardoublepage

\phantomsection
\addcontentsline{toc}{chapter}{Statement of Contributions}
\begin{center}\textbf{Statement of Contributions}\end{center}

\noindent
Yang Voelker provided the artwork for Figure~\ref{fig:architectures} as well as the brain illustration at the end of this front matter (the winning T-shirt design for the 2018 Telluride Neuromorphic Workshop).
Figure~\ref{fig:touch-network} was designed in collaboration with Ken E. Friedl.
Derivations from Lemma~\ref{lemma:coord-transform}, equation~\ref{eq:pulse-extended-double-exp}, and section~\ref{sec:mismatch}, extend work in collaboration with my co-authors, Kwabena Boahen and Terrence C. Stewart, both of whom provided important steps in accounting for pulse-extended double-exponential synapses with time-constant mismatch.
My advisor, Chris Eliasmith, helped author and edit several publications that have been adapted and extended throughout; references are included at the top of each applicable section.

\cleardoublepage

\phantomsection
\addcontentsline{toc}{chapter}{Abstract}
\begin{center}\textbf{Abstract}\end{center}

% The NEF is akin to a programming language for dynamical systems

\cleardoublepage

\phantomsection
\addcontentsline{toc}{chapter}{Acknowledgements}
\begin{center}\textbf{Acknowledgements}\end{center}

This thesis would not have been possible if not for my interactions with a few key people.
Most notably, my advisor, Chris Eliasmith, who gave me the freedom to swim into the deep end, and the structure, inspiration, counsel, and courage, to ask difficult, but relevant questions.
His active engagement and guidance has been instrumental throughout the development of this thesis.
Next, Terry Stewart, who programmed the Nengo backend for Braindrop, and has been an invaluable source of support when it comes to his technical prowess, breadth of knowledge, and intuition behind constructing functional spiking networks.
In my earlier days, Trevor Bekolay, Travis DeWolf, Daniel Rasmussen, and Xuan Choo, all played critical roles as academic mentors, by facilitating my exploration of the Neural Engineering Framework and Nengo.
Likewise, Eric Hunsberger, Jan Gosmann, Peter Duggins, and Andreas St{\"o}ckel, have each given feedback and provided foundational research that has deeply influenced my work.

Special mention must go out to Kwabena Boahen, who, throughout every one of our interactions, has indescribably imparted invaluable streams of insight in his critical and uncompromising quest for answers.
Many of his questions, contributions, and comments on papers, have inevitably influenced large portions of this thesis in a subtly-pervasive but positive manner.
This extends to memorable moments with all of his students, but especially Sam Fok, Alex Neckar, and Eric Kauderer-Abrams -- it has been a privilege to take part in these collaborations.

Briefly speaking with Mike Davies, hearing his eloquent characterization of neuromorphic computing, and having his genuine support in getting Loihi into the hands of researchers, has all been integral to propelling this work.
Wilten Nicola has provided a few influential communications on double-exponential synapses and FORCE learning, while Jamie Knight and Michael Hopkins have each provided important references involving SpiNNaker and high-dimensional sampling.
The words of many others have shaped this thesis throughout.

At the institutional level, I would like to acknowledge the support that I've received from
ABR (for Nengo and NengoLib development and experimentation with Loihi),
Intel (for access to Loihi),
Femtosense (for access to and support with Braindrop),
and the University of Manchester (for access to SpiNNaker).
Likewise, my acceptance into the neuromorphic workshops of Telluride, Capo Caccia, Intel's Neuromorphic Research Community~(INRC), and the Nengo Summer School, have all proven to be constructive experiences and motivational fuel.
The Centre for Theoretical Neuroscience~(CTN), Computational Neuroscience Research Group~(CNRG), David R. Cheriton School of Computer Science, and the University of Waterloo, have all contributed to providing a supportive academic environment over the past 11 years.  

\newpage

The experiments throughout this thesis make extensive use of Python, NumPy, SciPy, Jupyter, Matplotlib, Seaborn, Pandas, Nengo, PyStorm, Hyperopt, Scitkit-learn, Keras, and TensorFlow -- for this, we are incredibly grateful to its thousands of contributors.
A repeated mention is in order for Trevor, who provided this modified \LaTeX{} template and facilitated my contributions to the Nengo software ecosystem.
Likewise, Daniel, Eric, and Jan, for their help with Nengo development, provisioning of useful code, and feedback.

Working backwards in time, I would also like to thank Mani Fazeli for his many letters of support and words of wisdom at the beginning of graduate school, as well as Carol Norman for her recurring words of encouragement.
And the enriched education that I received from my teacher Jon Devlin, at Bell High School, was monumental in my decision to study computer science at the University of Waterloo.

At a personal level, Chris has been an outstanding role model for how to live a balanced and meaningful life, as an academic, parent, and human being, all in ways that I can only aspire towards.
Jason Batten-Carew has been the source of many extended philosophical discussions centered on the nature of reality, perception, stories, games, and cycles, all of which have undoubtedly rooted their way into my most basic pursuits for meaning.
I also thank my mom and dad for exciting my curiosity, and enabling my studies, throughout my entire life.
And finally I thank my wife, Yang Peng, for providing much-needed perspective and support throughout it all.

\cleardoublepage

\phantomsection
\addcontentsline{toc}{chapter}{Dedication}
\begin{center}
  \textbf{Dedication}

  To Yang -- the love of my life.
  And to Gunter -- for he is a dog.
\end{center}

\cleardoublepage

% https://tex.stackexchange.com/questions/65544/how-to-link-table-of-contents-in-thesis-pdf
\cleardoublepage
\phantomsection
\renewcommand\contentsname{Table of Contents}
%\pdfbookmark[section]{\contentsname}{Table of Contents}
\addcontentsline{toc}{chapter}{Table of Contents}
\tableofcontents
\cleardoublepage
\phantomsection

\iffalse
\addcontentsline{toc}{chapter}{List of Todos}
\listoftodos
\cleardoublepage
\phantomsection
\fi

\addcontentsline{toc}{chapter}{List of Tables}
\listoftables
\cleardoublepage
\phantomsection

\addcontentsline{toc}{chapter}{List of Figures}
\listoffigures
\cleardoublepage
\phantomsection

\addcontentsline{toc}{chapter}{List of Theorems}
\listoftheorems[ignoreall,show={theorem,lemma}]
\cleardoublepage
\phantomsection

\chapter*{Mathematical Conventions}
\addcontentsline{toc}{chapter}{Mathematical Conventions}
\label{typography}

This thesis follows standard mathematical conventions of using lower-case italic variables~($u(t)$) to represent scalar quantities, lower-case bold variables~($\V{x}(t)$) to represent vector quantities, upper-case italic variables~($A$) to represent matrices, and Greek letters~($\tau$, $\theta$) to represent quantities with a unit of time.
Bars are included above variables ($\bar{A}$) to indicate discretization in time.
Depending on context, $F(s)$ or $F(z)$ are used to represent complex values in the Laplace $s$-domain, or digital $z$-domain, respectively.

\cleardoublepage
\phantomsection

\begin{figure}
\centering
\vspace*{\fill}
\begingroup
\includegraphics[width=0.6\columnwidth]{yang-peng-telluride-back-2018.png}
\endgroup

\vspace{1em} 

\caption*{{\tiny \emph{``You would say to God, `How did you make the mountains?' and he would say... `Well, you're asking me to describe in words how I made the mountains, and there are no words which can do this. Words cannot tell you how I made the mountains any more than I can drink the ocean with a fork... It would take millions of years, and you would be bored with my description, long before I got through it... because I didn't create the mountains with words, I just did it. Like you open and close your hand. You know how you do this, but can you describe in words how you do it? Even a very good physiologist can't describe it in words. But you do it. You're conscious, aren't you. Don't you know how you manage to be conscious?''}~\textasciitilde~Alan Watts}}

\vspace*{\fill}
\end{figure}

\cleardoublepage
\phantomsection

% Change page numbering back to Arabic numerals
\pagenumbering{arabic}

% https://tex.stackexchange.com/questions/32208/footnote-runs-onto-second-page
\interfootnotelinepenalty=10000

\chapter{Introduction}

Computation cannot be physically realized without time.
Whether we consider the discrete switching of voltage states across transistors in a
digital computer,
or the continuous flow of charges across ion channels within the brain -- 
all physical systems, that we currently know of, make progress on their computations
by changing over time.
Digital computers time-multiplex their computations across multiple cores with energetically-expensive access to distant memory, while our brains consume only twenty watts of power by harnessing the dynamics of a biophysical substrate.
An overarching focus of this thesis is a computational paradigm
in between these two extremes: neuromorphic computing.
Here, researchers draw inspiration from how the
brain dynamically performs computations,
to engineer low-power digital and analog circuits that emulate its fundamental principles.

The emerging field of neuromorphic computing has already
witnessed many successes backed by large companies, small start-ups, and new research
programs being established at various institutions all around the world.
Applied Brain Research (ABR), Femtosense, Intel, IBM, and many others~\citep{marketreport2018}, are establishing themselves as key players in the field, while a collaborative landscape is being fostered across a diverse set of academic communities attempting to use such hardware including the Telluride Neuromorphic Cognition Engineering Workshop~\citep{cohen2001report}, Capo Caccia Cognitive Neuromorphic Engineering Workshop, Nengo Summer School~(ABR's ``Brain Camp''), and Intel's Neuromorphic Research Community~(INRC).

Despite this growing excitement, there exist many challenges when it comes to using
neuromorphic hardware to perform some desired task.
Generally speaking, neuromorphic programming forces us to rethink our traditional,
\citet{von1958}, understanding of computation, and redesign our algorithms to leverage distributed networks of noisy, heterogeneous units, that locally process their inputs and communicate via spikes -- also known as Spiking Neural Networks~(SNNs).
Biological nervous systems already do this quite successfully, and so we seek to borrow existing
principles and models from neuroscience, exploit their properties using tools from mathematics and engineering,
and use software methods to synthesize the resulting networks ``in silico''.
Moreover, the frameworks and methods that we use to accomplish this should be extensible---willing and able to incorporate increasingly-sophisticated levels of biological detail---so long as such details are computationally useful at a functional level, or shown to benefit some resource-power-accuracy trade-off.

To what end?
By virtue of emulating core principles of brain function, this realizes algorithms that consume far less power, and scale more readily to massively-parallelized computing architectures running in biological real-time.
But such algorithms must also be practically useful.
To flip this on its head: we require frameworks to describe what some hardware is doing at the lowest level,
and tools to systematically harness its capabilities to perform useful computations at the highest level.
Analagous to how a single compiler can translate the same C++ code
to dozens of physically distinct digital architectures, we strive to systematically describe
how the same dynamical computations can be flexibly mapped onto distinct spiking neuromorphic architectures.

The primary goals of this thesis are to explore the above statements, and in doing so elucidate a general recipe for implementing efficient dynamical computations on neuromorphic hardware, empower theorists and practitioners with deeper insights into subtle relationships between network configuration and function, and finally validate this approach with practical results that include a set of novel spiking algorithms deployed on state-of-the-art neuromorphic chips: Braindrop~\citep{braindrop2019} and Loihi~\citep{davies2018loihi}.
Ultimately, this unveils a class of computations where neuromorphic hardware excels, and further automates the process of engineering algorithms for such hardware.

We have already learned a great deal through the mutually interacting processes of bottom-up neuromorphic hardware design and top-down SNN construction.
Theoretical frameworks help guide design decisions, while hardware constraints demand relevant extensions to theoretical frameworks and software tools; both processes give-and-take in a dynamic feedback loop.
Our general approach consists of a software stack, and nonlinear dynamical systems framework,
for describing computations within a language of low-frequency vector spaces.
This approach is known as the Neural Engineering Framework~\citep[NEF;][]{eliasmith2003a}.
The NEF is implemented by the Python package, Nengo~\citep{bekolay2014}, which supports a variety of neuromorphic backends, and is extended by the methods
described in this thesis.
The ultimate goal of this approach is to develop a unified set of tools for constructing
artificially intelligent agents, that scale to the same kinds of tasks that
humans can perform, while running on neuromorphic hardware consuming several orders of magnitude less power than a traditional computer.
Significant progress has been made in this direction over the past two decades -- and although we still
have a long ways to go, this thesis aims to highlight a promising path for the neuromorphic community.

When confronted with the ambitious goal of artificially reproducing human intelligence, computer scientists will most often point towards
the popularity of deep learning~\citep{lecun2015deep} and its recent successes such as AlphaGo~\citep{gibney2016google}.
First and foremost, we seek to embrace the methods of deep learning, by incorporating its architectures and training methodoligies into our toolbox.
This is in fact the mandate of Nengo DL~\citep{rasmussen2018nengodl}, a Python package that combines Nengo and Tensorflow into a hybridized framework that automates backpropagation across spiking, non-spiking, and NEF-style recurrent neural networks -- broadly referred to as the class of Artificial Neural Networks~(ANNs).
Deep learning methods have found enormous success, due in part to the black-box nature of its training procedure, which requires a programmer to have virtually no knowledge of the underlying problem being solved.
However, we do not believe that deeper networks, better training heuristics, nor larger datasets, will be enough to achieve our aforementioned goal.
Some concerns for conventional deep learning approaches include: the amount of power required to scale traditional ANNs to the size of the human brain~\citep{furber2012build}, the difficulties tuning hyperparameters at scale~\citep{bergstra2015hyperopt}, and the abundance of adversarial attacks on networks trained via backpropagation~\citep{su2019one}.
More generally, the ``no free lunch theorem'' informs us that a black-box learning algorithm, on its own, will never be enough~\citep{wolpert1996lack}; inevitably, it becomes necessary to incorporate apriori knowledge into the structure of ANNs in a manner that biases the training procedure towards the problems being solved.
% Currently this prior knowledge is incorporated into deep learning through a variety of architectures, and the ability to tweak hyperparameters, but we can do better.

But even more fundamentally, the majority of ANNs in use today, including AlphaGo, are applied to \emph{static} inputs to produce \emph{static} outputs.
For a board game such as Go, this is forgivable, as the entire state of the game is encapsulated by the current state of the board itself, and so there is technically no need to consider any history in order to optimize across future states.\footnote{
Practically speaking, there could still be a benefit, for a sub-optimal player, to knowing the history of the game in order to better predict the opponent's strategy.}
Nevertheless, this points towards a thematic concern: time is not being treated as the dimension of computation in many of the most successful methods, which may limit their applicability to many real-world problems.
For instance, convolutional ANNs are most often deployed in video processing tasks by classifying objects in single frames,  %, without internal knowledge of how these objects might move or behave.
sampled at some arbitrary interval, and independently from one another.
Recurrently connected networks are an exception that we embrace, but they come with many challenges that we will discuss.
At large, most methods do not respect time as being a physical dimension that underpins the dynamics of the system's input, output, or processing units;
time is neither a variable in the neuron models, synapse models, nor any of the computational elements.
% Consequently, training these networks becomes an art-form than a science, 
% This pattern is beginning to reverse, with examples such as deep LSTMs and qRNNs for machine translation and speech synthesis.
% However, we argue that we can do even better, by arming ourselves with tools for better understanding the role of time, and being able to engineer specific dynamic computations into RNNs.

On the other hand, the human brain is a physical system that must continuously process information over time.
This is crucial for a cognitive agent embodied in a dynamic environment wherein all interactions depend on internally and externally evolving states.
The brain naturally learns from interacting with its environment over many different time-scales. 
For instance, our brains are not programmed to recognize static objects at discrete moments in time.
Rather, we have evolved to interact with objects in continuously changing environments, from which the perception of object recognition emerges.
From a very young age, children can intuitively infer the physics of moving objects, and learn how to interact with them to achieve desired outcomes.
This continues well through adulthood, up to the level of coordinating our thoughts and behaviours, sculpted by our cumulative experience with the world.
% Everything that we connect with helps to shape our brains to conceptualize these changes.
This occurs so naturally that we often take it for granted, but at a very basic basic level, our brains are necessarily constructing explicit and implicit dynamical models of reality, and applying these models
to perform behaviourally relevant computations in real-time.

Consider a real-world example such as driving a car. We must constantly integrate new information with our existing understanding of where we are, where we are headed, and what might happen along the way.
This requires not only an intuitive understanding of the physics of the car, but models of how other drivers behave on the road, with respect to changing traffic lights and road conditions, that play out in parallel while we flexibly plan and coordinate our actions.
Likewise, consider any actively engaging task or hobby that you enjoy.
I speculate this activity requires some amount of mentally or physically coordinated interaction with the world, on a similar level of dynamic complexity as in the driving example, such as: playing a video game, participating in a sport, engaging with some form of media, playing music, drawing, dancing, and so forth.
Besides such activities being meaningful, fun, and beautiful, they all share the common motif of recruiting a variety of systems in a dynamically coordinated fashion.
% Abstractly, we would like to conceptualize this as a compression hierarchy and decompression hierarchy whose two outside ends are coupled in time. Inside is ones own personal representation of the activity (see Fig. ???). At the core, the system strives to create the stable representation of an `illusion' that all of this takes place at once.
% If one is having fun, then this internal representation will try its best to encompass both ends.

To help appreciate how extraordinary such tasks are for the human brain, the mechanisms involved do not have a perfect memory of past inputs.
Each neuron locally processes some signal by continuously encoding it into spike trains.
These action potentials evoke post-synaptic currents~(PSCs) that decay exponentially on the order of milliseconds.
Memories emerge from these dynamics, as changes to activation patterns and synaptic weights reflect the histories of representations along a myriad of time-scales.
The time-scales over which these mechanisms interface with the world, and interact with each other, determines, and constrains, all of our perceptions, thoughts, and actions.
Remarkably, this is all accomplished by a collection of approximately one-hundred billion neurons, connected by over one-hundred trillion\footnote{
It is difficult to comprehend the sheer magnitude of this number, but one-hundred trillion is more than available estimates on the number of galaxies in the observable universe, trees on the earth, fish in the ocean, and digits of $\pi$ that have been computed -- all combined (at the time of writing).
}
synapses, processing these signals sparsely across space and time, while consuming only twenty watts~\citep{koch2014} -- about the same energy as a typical household CFL bulb.

This ``signal processing view'' of computation is fundamentally different from conventional views of computation targeted for conventional computers.
In the former view, time is intrinsic to the state of all physical ``devices'' performing the computation~(i.e.,~neurons and synapses), which are themselves necessarily tied to the timing and dynamics of the overall function as a consequence of actually implementing the function themselves.
However, in the conventional view, the time that a computation takes is an implementation detail that is mostly decoupled from the input-output relation.
As such, there is currently an important distinction between the computational approaches taken by the neuromorphic community and by those who program traditional digital computers.
This is not simply a matter of biological detail, but more crucially a matter of thinking about the system that is actually solving these tasks as a collection of signal processors whose dynamics are coupled to the task at hand by the physics of time.
If our aim is to build artificial systems that are as useful, robust, capable, and scalable as the human brain, then we should strive to understand computation in such terms.

The approach we take considers time as a fundamental starting point, and thus situates dynamics front-and-center; we bridge the time-scales of the entire system, down from the levels of spikes and post-synaptic currents, up to the levels of human reaction times and behaviours. 
We validate this model of computation by deploying these SNNs on low-power neuromorphic hardware designed to emulate key biological principles.
% There is no arithmetic logic unit, random-access memory, or instruction set.
This involves developing an extensible theoretical framework that links together the mechanisms of an SNN, its parameters, and some desired class of computations, while embracing all of the tools at our disposal.

%Another significant motivating factor for this work is the recent surge in neuromorphic computing architectures such as SpiNNaker~\citep{mundy2015efficient}, Neurogrid~\citep{choudhary2012silicon}, Brainstorm~\citep{brainstorm}, and IBM's TrueNorth~\citep{merolla2014million}.
%These platforms are massively-parallel low-power analog and/or digital systems that are designed specifically to simulate large-scale SNNs rivalling the complexity of the human brain.
%Despite the growing excitement surrounding these platforms, the community is struggling to fully unlock the computational power of this hardware.
%We believe this is primarily the result of peoples' inability to reconcile their discrete {\it von Neumann}-like understanding of conventional algorithms with the continuous brain-like processing of neuromorphic hardware.
%To take full advantage of such hardware, our `compilers' must account for the effects of simulated neuron models and synaptic models at the system level, and moreover exploit these effects in useful ways whenever possible.
%This has forced us to come full circle; once again we must understand computation in terms of the brain's processing units.

%The remainder of this thesis will be organized as follows.
%We will first survey existing methods of building RNNs (both spiking and non-spiking) and recent neuromorphic hardware architectures.
%We will then review the main ideas and tools that we use from control theory, linear systems theory, nonlinear dynamics, and systems design.
%Therefore, this thesis will theoretically explore ideas and to solve problems that arise from taking the above view of how the brain processes information. This will be supplemented by numerical simulations which validate theories, specific examples which demonstrate applicability, and connections to experimental studies which support or constrain our models when appropriate.
%Many of these problems are driven by current unpublished challenges within the neuromorphic community.
%A secondary goal is to demonstrate that these methods are practically useful by surpassing state-of-the-art results in machine learning, although we will not consider it a failure if this is not achieved within the remaining time span of the thesis, since there is still much ground to cover in theoretical foundations.
%Overall we aim to reveal new insights and make new methods widely available, that we hope will prove useful to engineers, theorists, and practitioners alike.


\chapter{Background}
\label{chapt:background}

\section{Recurrent Neural Networks}

\subsection{Traditional Approaches}

Backpropagation through time (unrolling)

Neural variants: Hunsburger, nengo\_dl, Maass, Hugh \& Sjenowski

\subsection{Structured Approaches}

LSTM, GRU, qRNN

\subsection{Unsupervised Approaches}

SORN, Numenta

\subsection{Reservoir Computing}

LSM, ESN

\subsection{Supervised State-Space Learning}

Duggins and activity-space extensions

FORCE, full-FORCE

Wilten Nicola

Aditya \& Gerstner, PES

\subsection{Balanced Networks}

Den\`eve, Memmesheimer, Slotine

\subsection{Others}

Friedemann Zen
Recurrent Information Optimization with Local, Metaplastic Synaptic Dynamics (ShiNung Ching)
System Level Approach (SLA) https://arxiv.org/abs/1701.05880


\section{Neural Engineering Framework}

\subsection{Principle 1}

Anectodally (e.g., personal conversation with Sophie Den\`eve, experience with Nengo, and averages in Spaun), we say 50-100 neurons per dimension achieves tolerable levels of noise.
For a fixed amount of noise in a represention implemented by an ensemble of $n$ neurons, the current best-known lower-bound predicts the dimensionality $d$ scales as $\Omega \left( n^{\frac{2}{3}} \right)$ for $n$ neurons~citep[][p.~60]{jgosmann2018}.
Although we don't know the upper-bound, we conjecture that it is $\bigoh \left( n \right)$.
That is, the dimensionality should not be able to scale faster than the neuron count.
If it could, then each neuron would represent $\omega(1)$ dimensions, which seems physically implausible given access to only $\bigoh (1)$ state variables.
This limitation could be broken by dendritic computation, for instance if each neuron had access to $\bigoh (n)$ variables distributed along the dendritic tree.

\subsection{Principle 2}

\subsection{Principle 3}


\section{Neuromorphic Computing}

Common goals.
Colocated memory and computation
Heterogeneity and sparsity in time
Heterogeneity and sparsity in space
Minimize spiking / synaptic events
Minimize power
Class of computations / descriptions
https://ieeexplore.ieee.org/document/7551407/ Christian Mayr's digital SRAM scaling 

\subsection{SpiNNaker}

1 and 2

\subsection{Neurogrid}

\subsection{Braindrop}

\subsection{Loihi}

\subsection{Others}

Spikey, BrainScaleS 1 and 2, Dynapse 1 and 2, TrueNorth, DeepSouth, COLAMN, ODIN, ROLLS, Giacomo and Eliasmith, Tripp, Wang and Tapson, STPU, Neuromemristive random projection networks (Dhireesha Kudithipudi)


\section{Dynamical Systems}

\subsection{Linear Time-Invariant Systems}

State-space representations, transfer functions, filters, convolution, properties, Pad\'e approximants and coordinate transformations

\subsection{Nonlinear Systems}

Linearization, Jacobians, signatures of chaos


\chapter{Principles}
\label{chapt:principles}

An ode to the three principles of the NEF.
These are themes that primarily expand upon the third principle.

\section{Single-Layer Equivalence}

The idea that any multi-layer network can be collapsed into a single recurrently connected layer.
In the NEF, this implies a certain structure with sparse encoders and sparse decoders.

\section{Heterogeneous Dynamical Primitives}

The idea that synapses, neurons, and even entire recurrent networks, can all be described in terms of their filtering / dynamical equations.
Of core importance is the notion that we want a rich basis of dynamics that is heterogeneous in time.

\section{Duality of Macroscopic and Microscopic Dynamics}

The corollary that, once all things are expressed in the same language, they become somewhat interchangeable.
Examples: (1) we can implement a lowpass filter as a synapse or as a network, (2) we can implement spiking LIF voltage dynamics as a single neuron or as a network.
But this also goes the other way (hence the duality).
Consider the architecture: $A \rightarrow (B \rightarrow B) \rightarrow A$, where $A$ and $B$ are neural ensembles.
We can think of $B \rightarrow B$ as a dynamical primitive, and then re-interpret the whole system as $A \rightarrow A$ with a much more complicated synapse model described by $B \rightarrow B$.
Like a ``virtual'' synapse, implemented by an entire sub-network.
The practical implication is that we can use the same theory and the same tools to understand and compose all of these things.
For example, if $B \rightarrow B$ is a delay network, we can then use what we know about amplifying axonal spike delays to make an even longer delay out of the entire system.
TODO: add to application section on delay networks.

\section{Continuous-Time versus Discrete-Time}

The idea that both continuous-time and discrete-time dynamics play important roles at different time-scales, both in terms of the dominant dynamics and in terms of the desired computation.
It is also important to consider both for neuromorphic hardware.

\section{Low-Frequency Representations}

The idea that the most amount of ``information'' is present within the low-frequency range of the signals, relative to the firing rates of neurons and time-constants of synapses.

\section{Low-Dimensional Representations}

Low-dimensional is relative to the number of neurons.
Anectodally (e.g., personal conversation with Sophie Den\`eve, experience with Nengo, and averages in Spaun), we say 50-100 neurons per dimension achieves tolerable levels of noise.
For a fixed amount of noise in a represention implemented by an ensemble of $n$ neurons, the current best-known lower-bound predicts the dimensionality $d$ scales as $\Omega \left( n^{\frac{2}{3}} \right)$ for $n$ neurons~\citep[][p.~60]{jgosmann2018}.
Although we don't know the upper-bound, we conjecture that it is $\bigoh \left( n \right)$.
That is, the dimensionality should not be able to scale faster than the neuron count.
If it could, then each neuron would represent $\omega(1)$ dimensions, which seems physically implausible given access to only $\bigoh (1)$ state variables.
This limitation could be broken by dendritic computation, for instance if each neuron had access to $\bigoh (n)$ variables distributed along the dendritic tree.


\end{document}
